% 第6節 推定と検定

% これはぜひ独立節で入れた方がいいです。
% しかも「理論」と「実装」をつなぐ節として非常に重要です。

% 構成は次がよいです。

% 6.1 尤度と最尤推定
% 尤度関数
% 対数尤度
% MLE
% 収束、数値最適化
% 6.2 識別可能性と正規化
% 効用は差だけが意味を持つ
% 基準代替、スケール
% パラメタが識別されるとは何か
% 6.3 推定結果の読み方
% 符号
% 限界効果
% Value of Time
% 弾力性
% WTP
% 6.4 検定

% 最低限これが欲しいです。

% t検定 / Wald検定
% 尤度比検定 (LR test)
% スコア検定
% モデル比較指標
% AIC
% BIC
% adjusted rho-square
% 6.5 予測性能と妥当性
% 的中率だけでは足りない
% in-sample fit と out-of-sample
% 行動解釈可能性
% transferability
% 6.6 SP/RPデータ、バイアス、内生性

% ここまで入るとかなり良いです。

% RP vs SP
% サンプル選択バイアス
% パネル相関
% 内生性
% measurement error